\documentclass[11pt,letterpaper]{article}
\usepackage{fontspec}
\usepackage{tocloft}
\usepackage{multicol}
\usepackage{nicefrac}
\usepackage{url}
\usepackage[left=1.3in,right=1.3in,top=1in,bottom=1in]{geometry}
\setmainfont[Scale=1.0,AutoFakeBold=1.5,AutoFakeSlant=0.3]{Doves Type}

% Prevent page breaks within list items
\widowpenalty=10000
\clubpenalty=10000
\interlinepenalty=500

\title{60-Minute Roast Chicken Stock (Instant Pot)}
\author{}
\date{}

\begin{document}

\maketitle
\thispagestyle{empty}

\begin{quote}\textit{A rotisserie chicken is stripped and its meat shredded fine, then pressure cooked \textit{40~minutes} with carrot, onion, and garlic---the extraction of six hours at a simmer, and clear because a sealed pot never boils. Makes about 5~cups of roast chicken stock.}\end{quote}

\section*{Ingredients}
\setlength{\columnsep}{20pt}
\noindent
    Chicken fat, skimmed from a previous batch, or neutral oil \dotfill 1~Tbsp. \\
    Carrot \dotfill \nicefrac{1}{4} \\
    Onion \dotfill \nicefrac{1}{4} \\
    Garlic cloves \dotfill 1 \\
    Rotisserie chicken \dotfill \nicefrac{1}{2} (about 1\nicefrac{1}{2}~lb.) \\
    Water \dotfill 1~quart \\
    Kosher salt \dotfill to taste (about \nicefrac{1}{2}~tsp. per quart) \\

\section*{Directions}

\noindent
Bring \textbf{water} to a boil ---
Slice \textbf{carrot} and \textbf{onion} thin ---
Smash \textbf{garlic} ---
Halve the \textbf{chicken} with kitchen shears and strip all of its meat, shredding it fine

\begin{enumerate}
    \item Heat the \textbf{chicken fat} in the pressure cooker on sauté over medium heat for \textit{1--2~minutes}. Add \textbf{carrot}, \textbf{onion}, and \textbf{garlic} and sweat for \textit{5--8~minutes}, until they soften and take a little colour at the edges; drop the heat if they start to brown all over.

    \item While the aromatics sweat, strip the meat from the \nicefrac{1}{2} \textbf{rotisserie chicken} and shred it finely with your fingers or two forks; small pieces shorten the flavour's path out into the water. Add the \textbf{meat}, \textbf{bones}, and \textbf{skin} to the pot.

    \item Pour in the boiling \textbf{water} and stir once. Seal the lid, set the valve to sealing, and cook on \textit{high pressure} for \textit{40~minutes}; the pot needs about \textit{10~minutes} to reach pressure before that timer starts.

    \item Switch the cooker off and let the pressure fall on its own, about \textit{15--20~minutes}. To hurry it, crack the valve just enough to let a thin thread of steam out, never wide open: the superheated liquid flashes to a boil and churns fat and fine particles back into the stock.

    \item Strain the \textbf{stock} through a fine-mesh strainer into a large bowl or 8-cup measuring cup and let it drain; do not press the solids, which hazes the stock. Discard them. You will have more than the quart of water you put in, about 5~cups, and it should taste of meat rather than bones.

    \item Chill the \textbf{stock}, lift the pale cap of \textbf{fat} off the surface, and keep it in a jar for your next batch. Taste before seasoning: the rotisserie bird has already salted the pot, so add \textbf{kosher salt} only if the stock tastes flat---about \nicefrac{1}{2}~tsp. to a quart finishes it.
\end{enumerate}

\newpage

{\scriptsize
\setlength{\columnsep}{20pt}
\setlength{\multicolsep}{6pt}
\begin{multicols}{2}
\setlength{\parindent}{0pt}
\setlength{\parskip}{2pt}
\setlength{\itemsep}{0pt}
\setlength{\parsep}{0pt}

\subsection*{Yield}
\begin{itemize}
    \item Makes about 5~cups (1\nicefrac{1}{4}~quarts), more than the water that went in
\end{itemize}

\subsection*{Equipment Required}
\begin{itemize}
    \item 6-quart electric pressure cooker (Instant Pot) or stovetop pressure cooker; a full 2-quart batch wants an 8-quart
    \item Kitchen shears, two forks, cutting board, and chef's knife
    \item Fine-mesh strainer, and a large bowl or 8-cup measuring cup
    \item Measuring cups and spoons
    \item Storage containers, plus a small jar for the skimmed \textbf{fat}
\end{itemize}

\subsection*{Mise en Place}
\begin{itemize}
    \item Slice \nicefrac{1}{4} \textbf{onion} and \nicefrac{1}{4} \textbf{carrot} thin and smash 1 clove \textbf{garlic}; keep them together on the board
    \item Halve the \textbf{chicken}, strip and shred the meat, and break the carcass into pieces that fit the pot
    \item Bring 1~quart \textbf{water} to a boil, and set out a jar for the \textbf{fat} you skim at the end
\end{itemize}

\subsection*{Ingredient Tips}
\begin{itemize}
    \item \textbf{Rotisserie chicken}: a supermarket bird is the point of the recipe---already roasted and seasoned, so it brings both the roast flavour and the salt. Any roast carcass and its leftover meat will do
    \item \textbf{Chicken fat}: keep the fat you skim from each batch and it carries roast flavour into the next one. Any neutral oil substitutes
    \item \textbf{Kosher salt}: the amount assumes Diamond Crystal; Morton is denser, so use about half. Taste first, since the bird has done most of the work
\end{itemize}

\subsection*{Preparation Tips}
\begin{itemize}
    \item The shredding is the technique. Small pieces shorten the diffusion path, which is why this tastes like meat in \textit{40~minutes} instead of after \textit{6~hours} at a simmer
    \item For chicken, stay between \textit{30} and \textit{90~minutes}. Beyond about \textit{3~hours} the bones break down and the stock turns metallic and bitter while the gelatin is chopped too fine to set
    \item Pressure holds the liquid below a boil, and that calm is what keeps the stock clear. A fast release undoes it
    \item A 6-quart cooker holds this half batch comfortably; double everything for a full batch and use an 8-quart, staying under the max-fill line
    \item An Instant Pot's high setting reaches about \textit{240°F}, a stovetop cooker at 15~psi nearer \textit{250°F}; the times here assume high pressure on either
    \item Judge it cold. The stock should set to a soft gel in the refrigerator, which is exactly what the consomm\'e asks for
\end{itemize}

\subsection*{Make Ahead \& Storage}
\begin{itemize}
    \item Keeps refrigerated \textit{4--5~days}, frozen up to \textit{3~months}; cool it uncovered and chill it fast
    \item Freeze in ice cube trays for portions---a cube or two drops straight into a pan sauce
    \item The skimmed \textbf{fat} keeps \textit{2~weeks} refrigerated or \textit{3~months} frozen
    \item Stock that stays liquid when cold was under-extracted or watered down; it still cooks well, it just will not body a sauce
    \item Best on the day it is made, while the roast notes are still bright
\end{itemize}

\subsection*{Serving Suggestions}
\begin{itemize}
    \item Pan sauces, soups, and risotto---anything that should taste of roast chicken
    \item Cook grains in it, or reduce it down with a splash of wine for a quick glaze
    \item A seasoned cup of it on its own is the flavour of a roast chicken without the bird
\end{itemize}

\subsection*{Source}
\begin{itemize}
    \item Adapted from Chris Young (Combustion Inc.; formerly The Fat Duck), via YouTube \textit{\$6 Michelin Stock in 60 Minutes (Costco Hack)}: \url{https://www.youtube.com/watch?v=3k20zFlbFfE}
\end{itemize}

\end{multicols}
}

\end{document}
